\section{Evidence}
\label{sec:evidence}

The tier model makes four falsifiable predictions. First, removing the Tier 1 component
from a working creative profile should produce outcomes worse than the bare baseline ---
not merely degraded but categorically worse. Second, Tier 2 components should be
individually necessary but not individually sufficient: each causes a below-threshold
failure when removed, yet none alone reaches threshold. Third, Tier 3 components should
be individually additive but jointly non-monotonic: adding them in the wrong order (Tier
3 before Tier 1) degrades performance. Fourth, personas that activate Tier 1 by identity
should outperform personas that activate Tier 3 by identity, by a margin that reflects
the full cost of missing the orientation layer. All four predictions are confirmed across
the reviewing and authoring data sets.

\subsection{The Reviewing-Side Tier Structure}

The eight-condition ablation study from \citeauthor{paper1} \cite{paper1} varied which
components of a structured expert reviewer profile were presented to an AI system, then
scored the resulting reviews on a 30-point rubric across five dimensions. The resulting
gradient runs C0 = 23.7 (bare baseline), C1 = 24.3 (credentials only), C2 = 23.3
(principles only), C3 = 26.7 (full principles), C4 = 26.7 (philosophy only), C5 = 21.7
(lens only), C6 = 25.3 (full profile), C7 = 22.3 (full profile plus all principles).

The first Tier 1 prediction is confirmed by C4 and C5. Philosophy alone (C4 = 26.7)
matches the full-principles condition (C3 = 26.7) despite containing far less information.
The orientation component replicates the performance of eleven principles in a single
sentence. More decisively, lens-only without philosophy (C5 = 21.7) falls 2.0 points
below the bare baseline (C0 = 23.7). Domain specificity without orientation is not merely
less helpful than orientation --- it actively degrades review quality, presumably by
focusing the system on format compliance at the expense of experiential relevance. The
lens is a Tier 3 tool: it sharpens defect detection within a category of concern, but
without the Tier 1 orientation that defines what the review is for, it surfaces formally
domain-appropriate observations that miss the solver experience entirely.

The Tier 2 prediction is confirmed by the structure of the non-redundant principles
identified in \citeauthor{paper1} \cite{paper1}: Verify the Last Mile, Blame the Player,
and the Computer Test. These three principles add unique diagnostic signal that the other
eight principles --- which load onto dimensions already captured by the philosophy ---
cannot supply. Their uniqueness is operational: each specifies a concrete test that can
be applied at any decision point in the review. The philosophy says ``evaluate from the
solver's experience''; these principles say ``did you check whether the solver can
actually complete the final extraction?'', ``did you consider whether the solver is at
fault before blaming the puzzle?'', ``did you test whether the mechanism is computable
without domain expertise?'' These are not restatements of the orientation; they are its
instantiation as checkable gates.

The non-monotonicity prediction is confirmed by comparing C6, C7, and C4. Full profile
(C6 = 25.3) underperforms philosophy alone (C4 = 26.7) by 1.4 points. Adding the review
lens and credentials to the philosophy dilutes rather than enriches. More strikingly, C7
(full profile plus all eleven principles = 22.3) substantially underperforms C6 (full
profile = 25.3). When Tier 2 non-redundant principles are bundled with Tier 3 redundant
ones and delivered without a stable Tier 1 foundation, the system cannot integrate the
additional constraints into a coherent evaluative stance. More context degrades
performance because additional Tier 3 content overwhelms rather than supplements an
orientation that has not been adequately activated. This is the tier model's central
non-monotonic prediction: adding Tier 3 material without a stable Tier 1 produces noise,
not signal.

\subsection{The Authoring-Side Tier Structure}

The trait decomposition study from \citeauthor{paper6} \cite{paper6} decomposed the
artist persona into six cognitive traits, measured quality under each trait in isolation
(Phase 1), under each single-trait removal from the full artist profile (Phase 2), and
under progressive reconstruction from the minimum sufficient set (Phase 3). The scoring
rubric was 45 points across nine dimensions.

Phase 1 isolates each trait's unaided contribution. Experience-first alone (T1 = 31.3)
substantially outperforms the bare baseline (A0 = 24.0) and falls below the full artist
(AA = 41.0) by approximately ten points, confirming that T1 is necessary but not alone
sufficient. Elegance drive alone (T4 = 26.0) falls \textit{below} the bare baseline ---
a result that is not a failure of T4 but an artifact of isolation: without the
orientation that tells the system what to optimize for, elegance becomes formal
minimalism that strips content rather than sharpening it. Rigor alone (T5 = 27.0) shows
the same pattern for the same reason. Neither Tier 2 trait can operate without Tier 1 to
direct it.

Phase 2 provides the cleanest test of tier membership. The leave-one-out drops from the
full artist (AA = 41.7) rank the traits precisely as the model predicts. Removing T1
produces a drop of $-$29.4 points (to 12.3), the only catastrophic failure in the
experiment. No other single-trait removal approaches this magnitude: T5 removal produces
$-$12.7, T4 removal $-$11.7, T2 removal $-$8.4, T6 removal $-$8.0, T3 removal $-$4.4.
The gap between T1's drop and the next largest drop (T5, $-$12.7) is itself 16.7 points.
Experience-first is not merely the most important trait; it occupies a qualitatively
different tier from all others.

Phase 3 tests the reconstruction curve: starting from zero, adding traits in Tier 1
then Tier 2 then Tier 3 order, and tracking where the 35-point quality threshold is
crossed. The baseline (R0 = 22.7) gives the starting point. Adding T1 alone (R1 = 28.3)
raises quality but does not cross threshold. Adding T4 (R2 = 30.3) and T5 (R3 = 35.7)
crosses the threshold at exactly the Tier 1 + Tier 2 boundary. This is the minimum
sufficient prompt: three traits, corresponding to OLE's three sentences. Subsequent
additions confirm the Tier 3 amplifying role: R4 = 37.3, R5 = 40.3, R6 = 42.3, each
adding between 2.0 and 3.0 points. The curve flattens as it approaches the full artist
ceiling (AA = 41.0 in Phase 1; 41.7 in Phase 2), consistent with diminishing returns in
the amplifying tier.

The threshold crossing at exactly R3 is the visual proof of the tier model's architecture.
Quality passes the publishable threshold when and only when both Tier 1 and both Tier 2
components are present. Adding any single Tier 3 component without that foundation
(demonstrated in Phase 1) does not cross threshold. Adding one Tier 2 component without
Tier 1 (also demonstrated in Phase 1) does not cross threshold. The model is not
additive across tiers; it is conditional.

\subsection{The Persona Natural Experiment}

The persona motivating study from \citeauthor{paper3} \cite{paper3} constitutes a natural
experiment in tier activation. Participants were given bare two-word persona labels ---
``you are an artist,'' ``you are a puzzle designer,'' ``you are a game designer,'' ``you
are a doctor,'' ``you are a programmer'' --- and scored on the same 45-point rubric. The
results, ranked by score, are: artist (AA) = 40.3, game designer (AG) = 32.3, programmer
= 24.0, puzzle designer (A0) = 24.7, doctor = 17.7.

The tier model explains this ordering without appeal to domain competence. The ``artist''
label activates Tier 1 by naming a practice whose defining criterion is the experience of
the viewer, listener, or solver. Artistic training is, functionally, training to inhabit
the receiver's experience before making any design decision --- which is exactly what the
orientation tier requires. The artist label also activates a partial Tier 2 through the
implicit elegance standard of artistic practice: ``does this need to be here?'' is a
question artists are trained to ask. The result (40.3) is near the ceiling for a single
bare label and approaches the carefully assembled full profile.

The puzzle designer label activates Tier 3 by naming a domain: puzzle design involves
acrostics, indexing, thematic assignment, answer extraction, and a set of genre
conventions. These are all Tier 3 elements. The label says nothing about the solver's
experience, nothing about removing unnecessary elements, nothing about verifying every
claim. The 15.6-point gap between artist (40.3) and puzzle designer (24.7) is the cost
of starting at Tier 3.

The game designer score (32.3) falls between the two, approximately 7.1 points above
puzzle designer and 8.0 points below artist. Game design training includes more
orientation toward the player experience than puzzle design training does --- playtesting
culture, player-experience vocabulary, the concept of ``game feel'' are all partial Tier
1 activations --- but game design is also defined by systems and mechanics that pull
toward Tier 3 operationalization. The partial Tier 1 activation explains approximately
half the gap relative to the artist.

The programmer (24.0) and doctor (17.7) scores confirm the model from the other
direction. Both professions activate professional precision --- formal correctness,
technical vocabulary, domain methodology --- which is Tier 3 domain-specificity with no
orientation layer. A programmer's quality criterion is whether the code runs correctly; a
doctor's is whether the clinical record is accurate. Neither criterion involves a theory
of what the receiver will experience. Applied to puzzle design, this precision produces
technically well-formed puzzles with no path for the solver --- formally correct and
experientially absent.

\subsection{The Designer Lens Comparison}

The three-designer comparison from \citeauthor{paper3} \cite{paper3} provides a within-tier
analysis of how different lens orientations in the reviewing role transfer to the
authoring role. Three experts --- Katz (operational lens), Young (perceptual lens),
Snyder (construction lens) --- were profiled in their reviewing roles (A5 conditions) and
then had their profiles augmented with authoring-specific guidance (A6 conditions). The
score deltas (A6 minus A5) are Katz $+$1.6, Young $+$3.7, Snyder $-$0.7.

These deltas are exactly what the tier model predicts. Katz's operational lens is a Tier
2 formulation: it specifies constraint-satisfaction tests applicable bidirectionally to
reviewing and authoring. When operational constraints are present, switching from the
reviewing to the authoring role requires only minor adaptation, because the constraints
--- ``does this mechanism have exactly one valid solution?'', ``does this step have a
recoverable error?'' --- apply in both directions. The $+$1.6 gain reflects the marginal
authoring-specific guidance added in A6. Operational lenses transfer cleanly across
roles.

Young's perceptual lens is Tier 3 in reviewing: it targets specific format-dependent
perceptual cues (contrast, visual grouping, typographic signaling) that reviewers use to
evaluate puzzle presentation. In the authoring role, these specific cues do not directly
transfer --- puzzle authors cannot replicate the exact perceptual grammar of a genre they
are not producing. But the $+$3.7 gain in A6 reveals what happens when Tier 3
format-specificity is blocked by a role boundary: Young's Tier 1 orientation (``how does
the solver perceive this?'') fills the gap by generating format-compatible substitutions.
The authoring profile produces novel perceptual solutions rather than replicating
reviewing heuristics. When Tier 3 specificity cannot transfer directly, Tier 1 orientation
activates creative substitution --- a result that cannot be explained by any additive or
averaging model of domain expertise.

Snyder's construction lens is already Tier 1 and Tier 2 in his reviewing profile. His
evaluation criteria focus on mechanism elegance (Tier 2: elegance gate) and solver-path
coherence (Tier 1: experience-first). These criteria are not role-bound: they apply
identically to reviewing a puzzle mechanism and constructing one. The near-zero delta
($-$0.7, within measurement noise) confirms that a profile with strong Tier 1 and Tier 2
components faces no role-crossing penalty. There is no gap to close because the
orientation and discipline are already authoring-compatible.

\begin{table}[t]
\centering
\caption{Cross-role symmetry: Tier assignment, evidence component, and score impact
across authoring and reviewing roles. All authoring scores are out of 45 (9 dimensions
$\times$ 5 points); all reviewing scores are out of 30 (5 dimensions $\times$ 6 points).}
\label{tab:symmetry}
\begin{tabular}{llrr}
\toprule
\textbf{Tier} & \textbf{Evidence Component} & \textbf{Reviewing $\Delta$} & \textbf{Authoring $\Delta$} \\
\midrule
Tier 1 & Orientation / Experience-first & $-$2.0 vs.\ baseline & $-$29.4 (L-T1) \\
Tier 2 & Operational discipline (non-redundant) & $+$3.0 (C3 vs.\ C0) & $-$12.2 avg.\ (L-T4, L-T5) \\
Tier 3 & Lens / Surprise, Visual, Systematic & $+$1.6 (C6 vs.\ C4) & $+$2.3 avg.\ (R4--R6 increments) \\
\midrule
Inverted & Tier 3 without Tier 1 & $-$2.0 (C5 vs.\ C0) & $-$17.0 (T4 alone vs.\ A0) \\
\bottomrule
\end{tabular}
\end{table}

The bottom row of Table~\ref{tab:symmetry} is the tier model's signature prediction, and
the one most difficult to explain by any flat expertise model: Tier 3 content without
Tier 1 activation produces performance \textit{below the bare baseline} in both roles.
The lens without philosophy (C5 = 21.7) underperforms the no-persona condition (C0 =
23.7). Elegance drive without experience-first (T4 alone = 26.0) underperforms the
no-trait condition (A0 = 24.0). Domain expertise without orientation actively degrades
creative quality.
